\begin{center}
\fcolorbox{ink}{pale}{\begin{minipage}{.92\linewidth}
\textbf{What is established, and what remains open.}
The mathematical contribution is a route from a tied-seed recursor plus a numerical equality test to full first-order intuitionistic PRA, through bounded simulation and finite quantifier-free function graphs. The written proof establishes induction preservation and both whole-domain return interpretations. The implementation contribution is a complete checker of the inherited calculus at 612 normalized bytes (604 with the inherited compact list predicate), versus the current 616-byte minimum.

These are different presentations. The small-checker figures are not assigned to the new arithmetic theory. The primary programs use 266 tokens, below the current 268-token minimum, but have 310/312 AST nodes. A separately qualified 639-byte ancillary variant improves the 300-node AST minimum to 298. The family envelope is therefore 604 bytes / 266 tokens / 298 AST nodes, achieved by different programs. An absolutely simplest axiomatization, an absolute size lower bound, and a single new checker below every old coordinate have not been established.
\end{minipage}}
\end{center}

\paragraph{Reading guide.}
Section~\ref{sec:new-basis} lists every primitive logical and nonlogical ingredient of the new theory. The arithmetic development explains the bounded initialization construction and its interpretation argument. The expressivity section identifies the exact standard term-function class and the role of finite case distinctions. Section~\ref{sec:checker} independently states the measured calculus, implementation change, audit, and limits. The appendices preserve the inherited checker bridge with attribution. All arithmetic metatheorems are written proofs; no proof-assistant verification is claimed.
